
\documentclass{article}
\usepackage{colortbl}
\usepackage{makecell}
\usepackage{multirow}
\usepackage{supertabular}

\begin{document}

\newcounter{utterance}

\centering \large Interaction Transcript for game `matchit\_ascii', experiment `similar\_grid\_1', episode 4 with Apertus{-}v1.5{-}8B.
\vspace{24pt}

{ \footnotesize  \setcounter{utterance}{1}
\setlength{\tabcolsep}{0pt}
\begin{supertabular}{c@{$\;$}|p{.15\linewidth}@{}p{.15\linewidth}p{.15\linewidth}p{.15\linewidth}p{.15\linewidth}p{.15\linewidth}}
    \# & $\;$A & \multicolumn{4}{c}{Game Master} & $\;\:$B\\
    \hline

    \theutterance \stepcounter{utterance}  
    & & \multicolumn{4}{p{0.6\linewidth}}{
        \cellcolor[rgb]{0.9,0.9,0.9}{
            \makecell[{{p{\linewidth}}}]{
                \texttt{\tiny{[P1$\langle$GM]}}
                \texttt{Συμμετέχετε σε ένα συνεργατικό παιχνίδι μαντεψιάς.} \\
\texttt{Ο στόχος είναι να βρείτε αν αυτό το πλέγμα και μια άλλη εικόνα που μόνο εγώ μπορώ να δω, είναι τα ίδια.} \\
\texttt{Παρακαλώ περιγράψτε πρώτα το πλέγμα σας. Στη συνέχεια, θα δώσω τη δική μου περιγραφή και μπορούμε να κάνουμε ο ένας στον άλλο ερωτήσεις σχετικά με τα πλέγματα για να καταλάβουμε αν είναι τα ίδια. Ξεκινήστε τώρα τη σύντομη περιγραφή του πλέγματός σας με "ΠΕΡΙΓΡΑΦΗ:" ακολουθούμενη από την περιγραφή. Μην προσθέσετε τίποτα άλλο. Αυτό είναι το πλέγμα σας:} \\
\texttt{X X ▢ X X} \\
\texttt{X ▢ ▢ ▢ X} \\
\texttt{▢ ▢ ▢ ▢ ▢} \\
\texttt{X X ▢ X X} \\
\texttt{X X ▢ X X} \\
            }
        }
    }
    & \\ \\

    \theutterance \stepcounter{utterance}  
    & \multicolumn{4}{p{0.6\linewidth}}{
        \cellcolor[rgb]{0.9,0.9,0.9}{
            \makecell[{{p{\linewidth}}}]{
                \texttt{\tiny{[P1$\rangle$GM]}}
                \texttt{ΠΕΡΙΓΡΑΦΗ: Οριζόντιες σειρές: Πρώτη σειρά: X, X, ▢, X, X Δεύτερη σειρά: X, ▢, ▢, ▢, X Τρίτη σειρά: ▢, ▢, ▢, ▢, ▢ Τέταρτη σειρά: X, X, ▢, X, X Πέμπτη σειρά: X, X, ▢, X, X} \\
            }
        }
    }
    & & \\ \\

    \theutterance \stepcounter{utterance}  
    & & & \multicolumn{2}{p{0.3\linewidth}}{
        \cellcolor[rgb]{0.9,0.9,0.9}{
            \makecell[{{p{\linewidth}}}]{
                \texttt{\tiny{[GM$|$GM]}}
                \texttt{continue} \\
            }
        }
    }
    & & \\ \\

    \theutterance \stepcounter{utterance}  
    & & \multicolumn{4}{p{0.6\linewidth}}{
        \cellcolor[rgb]{0.9,0.9,0.9}{
            \makecell[{{p{\linewidth}}}]{
                \texttt{\tiny{[GM$\rangle$P2]}}
                \texttt{Συμμετέχετε σε ένα συνεργατικό παιχνίδι μαντεψιάς.} \\
\texttt{Ο στόχος είναι να βρείτε αν αυτό το πλέγμα και μια άλλη εικόνα που μόνο εγώ μπορώ να δω, είναι τα ίδια.} \\
\texttt{Παρακαλώ περιγράψτε πρώτα το πλέγμα σας. Στη συνέχεια, θα δώσω τη δική μου περιγραφή και μπορούμε να κάνουμε ο ένας στον άλλο ερωτήσεις σχετικά με τα πλέγματα για να καταλάβουμε αν είναι τα ίδια. Ξεκινήστε τώρα τη σύντομη περιγραφή του πλέγματός σας με "ΠΕΡΙΓΡΑΦΗ:" ακολουθούμενη από την περιγραφή. Μην προσθέσετε τίποτα άλλο. Αυτό είναι το πλέγμα σας:} \\
\texttt{X X ▢ X X} \\
\texttt{X ▢ ▢ X X} \\
\texttt{▢ ▢ ▢ ▢ ▢} \\
\texttt{X ▢ ▢ X X} \\
\texttt{X X ▢ X X} \\
            }
        }
    }
    & \\ \\

    \theutterance \stepcounter{utterance}  
    & & & \multicolumn{4}{p{0.6\linewidth}}{
        \cellcolor[rgb]{0.9,0.9,0.9}{
            \makecell[{{p{\linewidth}}}]{
                \texttt{\tiny{[GM$\langle$P2]}}
                \texttt{ΠΕΡΙΓΡΑΦΗ: Το πλέγμα αποτελείται από 5 σειρές και 5 στήλες. Η πρώτη σειρά έχει τα στοιχεία: X, X, ▢, X, X. Η δεύτερη σειρά έχει τα στοιχεία: X, ▢, ▢, X, X. Η τρίτη σειρά έχει πέντε συνεχόμενα ▢. Η τέταρτη σειρά έχει τα στοιχεία: X, ▢, ▢, X, X. Η πέμπτη σειρά έχει τα στοιχεία: X, X, ▢, X, X.} \\
            }
        }
    }
     \\ \\

    \theutterance \stepcounter{utterance}  
    & & & \multicolumn{2}{p{0.3\linewidth}}{
        \cellcolor[rgb]{0.9,0.9,0.9}{
            \makecell[{{p{\linewidth}}}]{
                \texttt{\tiny{[GM$|$GM]}}
                \texttt{continue} \\
            }
        }
    }
    & & \\ \\

    \theutterance \stepcounter{utterance}  
    & & \multicolumn{4}{p{0.6\linewidth}}{
        \cellcolor[rgb]{0.9,0.9,0.9}{
            \makecell[{{p{\linewidth}}}]{
                \texttt{\tiny{[GM$\rangle$P2]}}
                \texttt{Αυτή είναι η ΠΕΡΙΓΡΑΦΗ μουΠΕΡΙΓΡΑΦΗ: Οριζόντιες σειρές: Πρώτη σειρά: X, X, ▢, X, X Δεύτερη σειρά: X, ▢, ▢, ▢, X Τρίτη σειρά: ▢, ▢, ▢, ▢, ▢ Τέταρτη σειρά: X, X, ▢, X, X Πέμπτη σειρά: X, X, ▢, X, X} \\
\texttt{Τώρα κάντε μια ερώτηση για να μάθετε νέες πτυχές του πλέγματός μου που μπορεί να διαφέρουν από το δικό σας πλέγμα. Ξεκινήστε με "ΕΡΩΤΗΣΗ:" και μην προσθέσετε τίποτα άλλο.} \\
            }
        }
    }
    & \\ \\

    \theutterance \stepcounter{utterance}  
    & & & \multicolumn{4}{p{0.6\linewidth}}{
        \cellcolor[rgb]{0.9,0.9,0.9}{
            \makecell[{{p{\linewidth}}}]{
                \texttt{\tiny{[GM$\langle$P2]}}
                \texttt{ΕΡΩΤΗΣΗ: Ποια στοιχεία βρίσκονται στη δεύτερη στήλη του πλέγματός σας;} \\
            }
        }
    }
     \\ \\

    \theutterance \stepcounter{utterance}  
    & & & \multicolumn{2}{p{0.3\linewidth}}{
        \cellcolor[rgb]{0.9,0.9,0.9}{
            \makecell[{{p{\linewidth}}}]{
                \texttt{\tiny{[GM$|$GM]}}
                \texttt{continue} \\
            }
        }
    }
    & & \\ \\

    \theutterance \stepcounter{utterance}  
    & & \multicolumn{4}{p{0.6\linewidth}}{
        \cellcolor[rgb]{0.9,0.9,0.9}{
            \makecell[{{p{\linewidth}}}]{
                \texttt{\tiny{[P1$\langle$GM]}}
                \texttt{Αυτή είναι η ΠΕΡΙΓΡΑΦΗ μουΠΕΡΙΓΡΑΦΗ: Το πλέγμα αποτελείται από 5 σειρές και 5 στήλες. Η πρώτη σειρά έχει τα στοιχεία: X, X, ▢, X, X. Η δεύτερη σειρά έχει τα στοιχεία: X, ▢, ▢, X, X. Η τρίτη σειρά έχει πέντε συνεχόμενα ▢. Η τέταρτη σειρά έχει τα στοιχεία: X, ▢, ▢, X, X. Η πέμπτη σειρά έχει τα στοιχεία: X, X, ▢, X, X.} \\
\texttt{ΕΡΩΤΗΣΗ: Ποια στοιχεία βρίσκονται στη δεύτερη στήλη του πλέγματός σας;Ξεκινήστε την απάντησή σας με "ΑΠΑΝΤΗΣΗ:" και μην προσθέσετε τίποτα άλλο.} \\
            }
        }
    }
    & \\ \\

    \theutterance \stepcounter{utterance}  
    & \multicolumn{4}{p{0.6\linewidth}}{
        \cellcolor[rgb]{0.9,0.9,0.9}{
            \makecell[{{p{\linewidth}}}]{
                \texttt{\tiny{[P1$\rangle$GM]}}
                \texttt{ΑΠΑΝΤΗΣΗ: X, ▢, ▢, ▢, X} \\
            }
        }
    }
    & & \\ \\

    \theutterance \stepcounter{utterance}  
    & & & \multicolumn{2}{p{0.3\linewidth}}{
        \cellcolor[rgb]{0.9,0.9,0.9}{
            \makecell[{{p{\linewidth}}}]{
                \texttt{\tiny{[GM$|$GM]}}
                \texttt{continue} \\
            }
        }
    }
    & & \\ \\

    \theutterance \stepcounter{utterance}  
    & & \multicolumn{4}{p{0.6\linewidth}}{
        \cellcolor[rgb]{0.9,0.9,0.9}{
            \makecell[{{p{\linewidth}}}]{
                \texttt{\tiny{[P1$\langle$GM]}}
                \texttt{Τώρα κάντε μια ερώτηση για να μάθετε νέες πτυχές του πλέγματός μου που μπορεί να διαφέρουν από το δικό σας πλέγμα. Ξεκινήστε με "ΕΡΩΤΗΣΗ:" και μην προσθέσετε τίποτα άλλο.} \\
            }
        }
    }
    & \\ \\

    \theutterance \stepcounter{utterance}  
    & \multicolumn{4}{p{0.6\linewidth}}{
        \cellcolor[rgb]{0.9,0.9,0.9}{
            \makecell[{{p{\linewidth}}}]{
                \texttt{\tiny{[P1$\rangle$GM]}}
                \texttt{ΕΡΩΤΗΣΗ: Υπάρχουν περιοχές στο πλέγμα σας που σχηματίζουν τετράγωνα ή ορθογώνια με όλα τα στοιχεία ▢;} \\
            }
        }
    }
    & & \\ \\

    \theutterance \stepcounter{utterance}  
    & & & \multicolumn{2}{p{0.3\linewidth}}{
        \cellcolor[rgb]{0.9,0.9,0.9}{
            \makecell[{{p{\linewidth}}}]{
                \texttt{\tiny{[GM$|$GM]}}
                \texttt{continue} \\
            }
        }
    }
    & & \\ \\

    \theutterance \stepcounter{utterance}  
    & & \multicolumn{4}{p{0.6\linewidth}}{
        \cellcolor[rgb]{0.9,0.9,0.9}{
            \makecell[{{p{\linewidth}}}]{
                \texttt{\tiny{[GM$\rangle$P2]}}
                \texttt{ΑΠΑΝΤΗΣΗ: X, ▢, ▢, ▢, X} \\
\texttt{ΕΡΩΤΗΣΗ: Υπάρχουν περιοχές στο πλέγμα σας που σχηματίζουν τετράγωνα ή ορθογώνια με όλα τα στοιχεία ▢;Ξεκινήστε την απάντησή σας με "ΑΠΑΝΤΗΣΗ:" και μην προσθέσετε τίποτα άλλο.} \\
            }
        }
    }
    & \\ \\

    \theutterance \stepcounter{utterance}  
    & & & \multicolumn{4}{p{0.6\linewidth}}{
        \cellcolor[rgb]{0.9,0.9,0.9}{
            \makecell[{{p{\linewidth}}}]{
                \texttt{\tiny{[GM$\langle$P2]}}
                \texttt{ΑΠΑΝΤΗΣΗ: Ναι, η τρίτη σειρά σχηματίζει ένα ορθογώνιο με όλα τα στοιχεία ▢.} \\
            }
        }
    }
     \\ \\

    \theutterance \stepcounter{utterance}  
    & & & \multicolumn{2}{p{0.3\linewidth}}{
        \cellcolor[rgb]{0.9,0.9,0.9}{
            \makecell[{{p{\linewidth}}}]{
                \texttt{\tiny{[GM$|$GM]}}
                \texttt{continue} \\
            }
        }
    }
    & & \\ \\

    \theutterance \stepcounter{utterance}  
    & & \multicolumn{4}{p{0.6\linewidth}}{
        \cellcolor[rgb]{0.9,0.9,0.9}{
            \makecell[{{p{\linewidth}}}]{
                \texttt{\tiny{[GM$\rangle$P2]}}
                \texttt{Τώρα κάντε μια ερώτηση για να μάθετε νέες πτυχές του πλέγματός μου που μπορεί να διαφέρουν από το δικό σας πλέγμα. Ξεκινήστε με "ΕΡΩΤΗΣΗ:" και μην προσθέσετε τίποτα άλλο.} \\
            }
        }
    }
    & \\ \\

    \theutterance \stepcounter{utterance}  
    & & & \multicolumn{4}{p{0.6\linewidth}}{
        \cellcolor[rgb]{0.9,0.9,0.9}{
            \makecell[{{p{\linewidth}}}]{
                \texttt{\tiny{[GM$\langle$P2]}}
                \texttt{ΕΡΩΤΗΣΗ: Ποια στοιχεία βρίσκονται στη τέταρτη στήλη του πλέγματός σας;} \\
            }
        }
    }
     \\ \\

    \theutterance \stepcounter{utterance}  
    & & & \multicolumn{2}{p{0.3\linewidth}}{
        \cellcolor[rgb]{0.9,0.9,0.9}{
            \makecell[{{p{\linewidth}}}]{
                \texttt{\tiny{[GM$|$GM]}}
                \texttt{continue} \\
            }
        }
    }
    & & \\ \\

    \theutterance \stepcounter{utterance}  
    & & \multicolumn{4}{p{0.6\linewidth}}{
        \cellcolor[rgb]{0.9,0.9,0.9}{
            \makecell[{{p{\linewidth}}}]{
                \texttt{\tiny{[P1$\langle$GM]}}
                \texttt{ΑΠΑΝΤΗΣΗ: Ναι, η τρίτη σειρά σχηματίζει ένα ορθογώνιο με όλα τα στοιχεία ▢.} \\
\texttt{ΕΡΩΤΗΣΗ: Ποια στοιχεία βρίσκονται στη τέταρτη στήλη του πλέγματός σας;Ξεκινήστε την απάντησή σας με "ΑΠΑΝΤΗΣΗ:" και μην προσθέσετε τίποτα άλλο.} \\
            }
        }
    }
    & \\ \\

    \theutterance \stepcounter{utterance}  
    & \multicolumn{4}{p{0.6\linewidth}}{
        \cellcolor[rgb]{0.9,0.9,0.9}{
            \makecell[{{p{\linewidth}}}]{
                \texttt{\tiny{[P1$\rangle$GM]}}
                \texttt{ΑΠΑΝΤΗΣΗ: X, X, ▢, X, X} \\
            }
        }
    }
    & & \\ \\

    \theutterance \stepcounter{utterance}  
    & & & \multicolumn{2}{p{0.3\linewidth}}{
        \cellcolor[rgb]{0.9,0.9,0.9}{
            \makecell[{{p{\linewidth}}}]{
                \texttt{\tiny{[GM$|$GM]}}
                \texttt{continue} \\
            }
        }
    }
    & & \\ \\

    \theutterance \stepcounter{utterance}  
    & & \multicolumn{4}{p{0.6\linewidth}}{
        \cellcolor[rgb]{0.9,0.9,0.9}{
            \makecell[{{p{\linewidth}}}]{
                \texttt{\tiny{[P1$\langle$GM]}}
                \texttt{Τώρα κάντε μια ερώτηση για να μάθετε νέες πτυχές του πλέγματός μου που μπορεί να διαφέρουν από το δικό σας πλέγμα. Ξεκινήστε με "ΕΡΩΤΗΣΗ:" και μην προσθέσετε τίποτα άλλο.} \\
            }
        }
    }
    & \\ \\

    \theutterance \stepcounter{utterance}  
    & \multicolumn{4}{p{0.6\linewidth}}{
        \cellcolor[rgb]{0.9,0.9,0.9}{
            \makecell[{{p{\linewidth}}}]{
                \texttt{\tiny{[P1$\rangle$GM]}}
                \texttt{ΕΡΩΤΗΣΗ: Υπάρχουν στοιχεία X που βρίσκονται ακριβώς δίπλα το ένα δίπλα στο άλλο σε οριζόντια σειρά;} \\
            }
        }
    }
    & & \\ \\

    \theutterance \stepcounter{utterance}  
    & & & \multicolumn{2}{p{0.3\linewidth}}{
        \cellcolor[rgb]{0.9,0.9,0.9}{
            \makecell[{{p{\linewidth}}}]{
                \texttt{\tiny{[GM$|$GM]}}
                \texttt{continue} \\
            }
        }
    }
    & & \\ \\

    \theutterance \stepcounter{utterance}  
    & & \multicolumn{4}{p{0.6\linewidth}}{
        \cellcolor[rgb]{0.9,0.9,0.9}{
            \makecell[{{p{\linewidth}}}]{
                \texttt{\tiny{[GM$\rangle$P2]}}
                \texttt{ΑΠΑΝΤΗΣΗ: X, X, ▢, X, X} \\
\texttt{ΕΡΩΤΗΣΗ: Υπάρχουν στοιχεία X που βρίσκονται ακριβώς δίπλα το ένα δίπλα στο άλλο σε οριζόντια σειρά;Ξεκινήστε την απάντησή σας με "ΑΠΑΝΤΗΣΗ:" και μην προσθέσετε τίποτα άλλο.} \\
            }
        }
    }
    & \\ \\

    \theutterance \stepcounter{utterance}  
    & & & \multicolumn{4}{p{0.6\linewidth}}{
        \cellcolor[rgb]{0.9,0.9,0.9}{
            \makecell[{{p{\linewidth}}}]{
                \texttt{\tiny{[GM$\langle$P2]}}
                \texttt{ΑΠΑΝΤΗΣΗ: Ναι, υπάρχουν στοιχεία X που βρίσκονται ακριβώς δίπλα το ένα δίπλα στο άλλο σε οριζόντια σειρά, όπως στη πρώτη, τέταρτη και πέμπτη σειρά.} \\
            }
        }
    }
     \\ \\

    \theutterance \stepcounter{utterance}  
    & & & \multicolumn{2}{p{0.3\linewidth}}{
        \cellcolor[rgb]{0.9,0.9,0.9}{
            \makecell[{{p{\linewidth}}}]{
                \texttt{\tiny{[GM$|$GM]}}
                \texttt{continue} \\
            }
        }
    }
    & & \\ \\

    \theutterance \stepcounter{utterance}  
    & & \multicolumn{4}{p{0.6\linewidth}}{
        \cellcolor[rgb]{0.9,0.9,0.9}{
            \makecell[{{p{\linewidth}}}]{
                \texttt{\tiny{[GM$\rangle$P2]}}
                \texttt{Τώρα κάντε μια ερώτηση για να μάθετε νέες πτυχές του πλέγματός μου που μπορεί να διαφέρουν από το δικό σας πλέγμα. Ξεκινήστε με "ΕΡΩΤΗΣΗ:" και μην προσθέσετε τίποτα άλλο.} \\
            }
        }
    }
    & \\ \\

    \theutterance \stepcounter{utterance}  
    & & & \multicolumn{4}{p{0.6\linewidth}}{
        \cellcolor[rgb]{0.9,0.9,0.9}{
            \makecell[{{p{\linewidth}}}]{
                \texttt{\tiny{[GM$\langle$P2]}}
                \texttt{ΕΡΩΤΗΣΗ: Υπάρχουν περιοχές στο πλέγμα σας που περιέχουν μόνο ένα X και τρία ▢ σε κάθε σειρά ή στήλη;} \\
            }
        }
    }
     \\ \\

    \theutterance \stepcounter{utterance}  
    & & & \multicolumn{2}{p{0.3\linewidth}}{
        \cellcolor[rgb]{0.9,0.9,0.9}{
            \makecell[{{p{\linewidth}}}]{
                \texttt{\tiny{[GM$|$GM]}}
                \texttt{continue} \\
            }
        }
    }
    & & \\ \\

    \theutterance \stepcounter{utterance}  
    & & \multicolumn{4}{p{0.6\linewidth}}{
        \cellcolor[rgb]{0.9,0.9,0.9}{
            \makecell[{{p{\linewidth}}}]{
                \texttt{\tiny{[P1$\langle$GM]}}
                \texttt{ΑΠΑΝΤΗΣΗ: Ναι, υπάρχουν στοιχεία X που βρίσκονται ακριβώς δίπλα το ένα δίπλα στο άλλο σε οριζόντια σειρά, όπως στη πρώτη, τέταρτη και πέμπτη σειρά.} \\
\texttt{ΕΡΩΤΗΣΗ: Υπάρχουν περιοχές στο πλέγμα σας που περιέχουν μόνο ένα X και τρία ▢ σε κάθε σειρά ή στήλη;Ξεκινήστε την απάντησή σας με "ΑΠΑΝΤΗΣΗ:" και μην προσθέσετε τίποτα άλλο.} \\
            }
        }
    }
    & \\ \\

    \theutterance \stepcounter{utterance}  
    & \multicolumn{4}{p{0.6\linewidth}}{
        \cellcolor[rgb]{0.9,0.9,0.9}{
            \makecell[{{p{\linewidth}}}]{
                \texttt{\tiny{[P1$\rangle$GM]}}
                \texttt{ΑΠΑΝΤΗΣΗ: Όχι, δεν υπάρχουν περιοχές στο πλέγμα σας που περιέχουν μόνο ένα X και τρία ▢ σε κάθε σειρά ή στήλη.} \\
            }
        }
    }
    & & \\ \\

    \theutterance \stepcounter{utterance}  
    & & & \multicolumn{2}{p{0.3\linewidth}}{
        \cellcolor[rgb]{0.9,0.9,0.9}{
            \makecell[{{p{\linewidth}}}]{
                \texttt{\tiny{[GM$|$GM]}}
                \texttt{continue} \\
            }
        }
    }
    & & \\ \\

    \theutterance \stepcounter{utterance}  
    & & \multicolumn{4}{p{0.6\linewidth}}{
        \cellcolor[rgb]{0.9,0.9,0.9}{
            \makecell[{{p{\linewidth}}}]{
                \texttt{\tiny{[P1$\langle$GM]}}
                \texttt{Τώρα κάντε μια ερώτηση για να μάθετε νέες πτυχές του πλέγματός μου που μπορεί να διαφέρουν από το δικό σας πλέγμα. Ξεκινήστε με "ΕΡΩΤΗΣΗ:" και μην προσθέσετε τίποτα άλλο.} \\
            }
        }
    }
    & \\ \\

    \theutterance \stepcounter{utterance}  
    & \multicolumn{4}{p{0.6\linewidth}}{
        \cellcolor[rgb]{0.9,0.9,0.9}{
            \makecell[{{p{\linewidth}}}]{
                \texttt{\tiny{[P1$\rangle$GM]}}
                \texttt{ΕΡΩΤΗΣΗ: Το πλέγμα σας περιέχει συμμετρία σε οποιαδήποτε από τις σειρές ή τις στήλες;} \\
            }
        }
    }
    & & \\ \\

    \theutterance \stepcounter{utterance}  
    & & & \multicolumn{2}{p{0.3\linewidth}}{
        \cellcolor[rgb]{0.9,0.9,0.9}{
            \makecell[{{p{\linewidth}}}]{
                \texttt{\tiny{[GM$|$GM]}}
                \texttt{continue} \\
            }
        }
    }
    & & \\ \\

    \theutterance \stepcounter{utterance}  
    & & \multicolumn{4}{p{0.6\linewidth}}{
        \cellcolor[rgb]{0.9,0.9,0.9}{
            \makecell[{{p{\linewidth}}}]{
                \texttt{\tiny{[GM$\rangle$P2]}}
                \texttt{ΑΠΑΝΤΗΣΗ: Όχι, δεν υπάρχουν περιοχές στο πλέγμα σας που περιέχουν μόνο ένα X και τρία ▢ σε κάθε σειρά ή στήλη.} \\
\texttt{ΕΡΩΤΗΣΗ: Το πλέγμα σας περιέχει συμμετρία σε οποιαδήποτε από τις σειρές ή τις στήλες;Ξεκινήστε την απάντησή σας με "ΑΠΑΝΤΗΣΗ:" και μην προσθέσετε τίποτα άλλο.} \\
            }
        }
    }
    & \\ \\

    \theutterance \stepcounter{utterance}  
    & & & \multicolumn{4}{p{0.6\linewidth}}{
        \cellcolor[rgb]{0.9,0.9,0.9}{
            \makecell[{{p{\linewidth}}}]{
                \texttt{\tiny{[GM$\langle$P2]}}
                \texttt{ΑΠΑΝΤΗΣΗ: Η πρώτη και η πέμπτη σειρά είναι συμμετρικές μεταξύ τους. Η τέταρτη και η πρώτη σειρά είναι επίσης συμμετρικές.} \\
            }
        }
    }
     \\ \\

    \theutterance \stepcounter{utterance}  
    & & & \multicolumn{2}{p{0.3\linewidth}}{
        \cellcolor[rgb]{0.9,0.9,0.9}{
            \makecell[{{p{\linewidth}}}]{
                \texttt{\tiny{[GM$|$GM]}}
                \texttt{continue} \\
            }
        }
    }
    & & \\ \\

    \theutterance \stepcounter{utterance}  
    & & \multicolumn{4}{p{0.6\linewidth}}{
        \cellcolor[rgb]{0.9,0.9,0.9}{
            \makecell[{{p{\linewidth}}}]{
                \texttt{\tiny{[GM$\rangle$P2]}}
                \texttt{Τώρα πάρε μια απόφαση. Τι πιστεύεις: είναι το δικό σου πλέγμα και το άλλο πλέγμα που περιγράφεται το ίδιο; Γράψε "ΑΠΟΦΑΣΗ: ίδιο πλέγμα" αν πιστεύεις ότι είναι το ίδιο πλέγμα ή "ΑΠΟΦΑΣΗ: διαφορετικά πλέγματα" αν πιστεύεις ότι είναι διαφορετικά πλέγματα. Μην προσθέσεις τίποτα άλλο.} \\
            }
        }
    }
    & \\ \\

    \theutterance \stepcounter{utterance}  
    & & & \multicolumn{4}{p{0.6\linewidth}}{
        \cellcolor[rgb]{0.9,0.9,0.9}{
            \makecell[{{p{\linewidth}}}]{
                \texttt{\tiny{[GM$\langle$P2]}}
                \texttt{ΑΠΟΦΑΣΗ: διαφορετικά πλέγματα} \\
            }
        }
    }
     \\ \\

    \theutterance \stepcounter{utterance}  
    & & & \multicolumn{2}{p{0.3\linewidth}}{
        \cellcolor[rgb]{0.9,0.9,0.9}{
            \makecell[{{p{\linewidth}}}]{
                \texttt{\tiny{[GM$|$GM]}}
                \texttt{continue} \\
            }
        }
    }
    & & \\ \\

    \theutterance \stepcounter{utterance}  
    & & & \multicolumn{2}{p{0.3\linewidth}}{
        \cellcolor[rgb]{0.9,0.9,0.9}{
            \makecell[{{p{\linewidth}}}]{
                \texttt{\tiny{[GM$|$GM]}}
                \texttt{success} \\
            }
        }
    }
    & & \\ \\

    \theutterance \stepcounter{utterance}  
    & & \multicolumn{4}{p{0.6\linewidth}}{
        \cellcolor[rgb]{0.9,0.9,0.9}{
            \makecell[{{p{\linewidth}}}]{
                \texttt{\tiny{[P1$\langle$GM]}}
                \texttt{ΑΠΑΝΤΗΣΗ: Η πρώτη και η πέμπτη σειρά είναι συμμετρικές μεταξύ τους. Η τέταρτη και η πρώτη σειρά είναι επίσης συμμετρικές.} \\
\texttt{Τώρα πάρε μια απόφαση. Τι πιστεύεις: είναι το δικό σου πλέγμα και το άλλο πλέγμα που περιγράφεται το ίδιο; Γράψε "ΑΠΟΦΑΣΗ: ίδιο πλέγμα" αν πιστεύεις ότι είναι το ίδιο πλέγμα ή "ΑΠΟΦΑΣΗ: διαφορετικά πλέγματα" αν πιστεύεις ότι είναι διαφορετικά πλέγματα. Μην προσθέσεις τίποτα άλλο.} \\
            }
        }
    }
    & \\ \\

    \theutterance \stepcounter{utterance}  
    & \multicolumn{4}{p{0.6\linewidth}}{
        \cellcolor[rgb]{0.9,0.9,0.9}{
            \makecell[{{p{\linewidth}}}]{
                \texttt{\tiny{[P1$\rangle$GM]}}
                \texttt{ΑΠΟΦΑΣΗ: ίδιο πλέγμα} \\
            }
        }
    }
    & & \\ \\

    \theutterance \stepcounter{utterance}  
    & & & \multicolumn{2}{p{0.3\linewidth}}{
        \cellcolor[rgb]{0.9,0.9,0.9}{
            \makecell[{{p{\linewidth}}}]{
                \texttt{\tiny{[GM$|$GM]}}
                \texttt{continue} \\
            }
        }
    }
    & & \\ \\

    \theutterance \stepcounter{utterance}  
    & & & \multicolumn{2}{p{0.3\linewidth}}{
        \cellcolor[rgb]{0.9,0.9,0.9}{
            \makecell[{{p{\linewidth}}}]{
                \texttt{\tiny{[GM$|$GM]}}
                \texttt{loss} \\
            }
        }
    }
    & & \\ \\

\end{supertabular}
}

\end{document}
